\begin{table}[!htbp]
\centering
\caption{School value added and second-academic-year persistence}
\label{tab:va-higher-ed-persistence-cross-va}
\begin{threeparttable}
\begin{tabular}{lccc}
\toprule
 & Higher-ed enrollment & High-premium field & High-premium institution \\
\midrule
Higher-ed enrollment VA & 0.244* & 0.055 & -0.039 \\
 & (0.143) & (0.057) & (0.058) \\
High-premium field VA & 0.031 & 0.032 & 0.037 \\
 & (0.080) & (0.032) & (0.035) \\
High-premium institution VA & 0.127 & -0.038 & 0.323*** \\
 & (0.162) & (0.068) & (0.107) \\
\midrule
Math VA & -0.094 & -0.003 & 0.033 \\
 & (0.058) & (0.024) & (0.027) \\
Verbal VA & -0.142** & -0.036 & 0.021 \\
 & (0.064) & (0.026) & (0.029) \\
Admission-exam taking VA & -0.011 & -0.019 & 0.007 \\
 & (0.064) & (0.025) & (0.028) \\
Projected program income VA & 0.450*** & 0.048 & 0.138** \\
 & (0.134) & (0.055) & (0.069) \\
Fin. aid app. VA & 0.093 & 0.014 & 0.000 \\
 & (0.136) & (0.054) & (0.053) \\
\bottomrule
\end{tabular}
\begin{tablenotes}[flushleft]
\footnotesize
\item Notes: Columns report persistence outcomes requiring entry into the relevant category and continued participation in the second academic year. Rows identify the EB observational school value-added measure used as the endogenous school-quality index. In each cell, attended-school VA is instrumented with first-round offered-school VA, controlling for the DA-probability expected value of the same row VA, cohort, grade-4 math and verbal scores, gender, and age. The sample contains timely SAE applicants from the 2018--2019 cohorts with nondegenerate assignment risk; admission-exam taking is not a sample restriction. Heteroskedasticity-robust standard errors are in parentheses. $^{*}p<0.10$, $^{**}p<0.05$, $^{***}p<0.01$.
\end{tablenotes}
\end{threeparttable}
\end{table}
